\chapter{Desenvolvimento do Trabalho} \label{cap:desenvolvimento}

\section{Visão geral do desenvolvimento}

O FLIKE foi desenvolvido como um sistema integrado de gestão de autorizações e controle de acesso físico. Sua implementação combina uma aplicação web, uma API com persistência relacional, uma credencial temporária transportada por QR Code e um dispositivo embarcado capaz de tomar a decisão de abertura junto à fechadura. Essa decomposição permite separar a concessão administrativa do acesso da apresentação da credencial no espaço físico, mantendo entre as duas etapas um formato comum de dados.

A parte conectada compreende cadastro, autenticação, administração da estrutura institucional, solicitação de acesso, aprovação ou rejeição e emissão da credencial. Nela, o frontend organiza a interação com usuários e responsáveis, enquanto a API aplica as regras do domínio e persiste contas, recursos, solicitações, trancas e credenciais. Depois da aprovação, o usuário recupera o payload emitido e a aplicação o representa como QR Code.

A parte local começa quando esse QR Code é apresentado à câmera da ESP32-CAM. O dispositivo recupera o payload, verifica sua estrutura, seu destino, sua janela temporal e seu código de autenticação e, quando a decisão é positiva, produz o sinal lógico enviado ao circuito de acionamento. Essa decisão não exige consulta à API no momento da leitura. A expressão ``operação local'' empregada neste capítulo refere-se a essa etapa; a obtenção inicial da credencial continua dependente da aplicação conectada.

As seções seguintes descrevem primeiro a arquitetura e o modelo de dados e depois detalham aplicação, credencial, dispositivo físico, segurança e avaliação. A descrição combina o código preservado nos repositórios, a documentação técnica e as demonstrações confirmadas pela equipe. A existência de código é tratada como evidência de implementação, enquanto o atendimento dos requisitos será confrontado posteriormente com testes e resultados.

\section{Arquitetura do FLIKE}

A arquitetura organiza as responsabilidades em torno do percurso da autorização: o servidor registra a decisão administrativa e emite uma credencial; o navegador a transforma em representação visual; e o dispositivo físico interpreta essa representação e decide sobre o acionamento. A Figura~\ref{fig:arquitetura-flike} sintetiza esse percurso e destaca a fronteira entre as operações conectadas e a validação local.

\subsection{Componentes e responsabilidades}

O usuário interage com o frontend executado no navegador. A interface envia requisições HTTP em JSON à API FastAPI, que reúne as regras de autenticação, propriedade institucional, solicitação, aprovação e emissão. A API acessa o banco MySQL para consultar e persistir o estado do domínio. Quando uma credencial é consultada, seu payload binário é representado em hexadecimal na resposta JSON; o frontend reconstrói os bytes e gera o QR Code localmente.

No espaço físico, o QR Code funciona como canal óptico entre a aplicação e a ESP32-CAM. A câmera captura o símbolo, e o firmware decodifica e avalia a credencial com o identificador, o relógio e o segredo provisionados no dispositivo. Uma autorização positiva produz um sinal lógico \texttt{HIGH}. O estágio com transistor e relé adapta esse comando ao circuito da fechadura elétrica. O Quadro~\ref{qua:componentes-flike} resume as responsabilidades sem antecipar os detalhes das seções dedicadas a cada subsistema.

\begin{quadro}[htb]
    \caption{Componentes e responsabilidades na arquitetura do FLIKE}
    \label{qua:componentes-flike}
    \small
    \begin{tabular}{|p{0.25\textwidth}|p{0.40\textwidth}|p{0.25\textwidth}|}
        \hline
        \textbf{Componente} & \textbf{Responsabilidade principal} & \textbf{Fronteira relevante} \\ \hline
        Frontend Next.js/React & Apresentar os fluxos web, comunicar-se com a API e renderizar a credencial como QR Code. & Depende da API para obter e atualizar os dados do sistema. \\ \hline
        API FastAPI & Autenticar contas, aplicar regras do domínio, decidir operações administrativas e emitir credenciais. & Não participa da decisão física no instante da leitura. \\ \hline
        Banco MySQL & Persistir usuários, estrutura institucional, trancas, solicitações, credenciais e registros previstos. & Representa estado lógico; não comprova por si só um evento físico. \\ \hline
        QR Code & Transportar visualmente o payload binário já emitido até a câmera. & Pode ser copiado; não executa autenticação por si próprio. \\ \hline
        ESP32-CAM e firmware & Capturar o símbolo, recuperar o payload, decidir localmente e emitir o comando lógico. & Opera com dados, relógio e segredo disponíveis no dispositivo. \\ \hline
        Circuito e fechadura & Adaptar o sinal lógico e realizar o destravamento elétrico. & O comando não comprova que a porta abriu ou que uma pessoa entrou. \\ \hline
    \end{tabular}
    \fonte{Elaborado pelos autores.}
\end{quadro}

\begin{figure}[htb]
    \centering
    \includegraphics[width=\textwidth]{arquitetura_flike.pdf}
    \caption{Arquitetura geral e fronteira de operação do FLIKE}
    \label{fig:arquitetura-flike}
    \fonte{Elaborado pelos autores.}
\end{figure}

\subsection{Fronteira entre operação conectada e operação local}

Na região conectada da Figura~\ref{fig:arquitetura-flike}, o navegador, a API e o banco mantêm o estado administrativo. O usuário precisa acessar essa região para cadastrar-se, autenticar-se, solicitar uma credencial, acompanhar a decisão do responsável e recuperar o QR Code emitido. O FLIKE não implementa um aplicativo móvel nativo nem oferece gestão web independente do servidor.

A fronteira é atravessada pelo próprio payload apresentado visualmente. Depois que a credencial está disponível, a ESP32-CAM não precisa consultar o banco nem a API para decidir sobre a abertura: ela verifica localmente as informações autenticadas. Essa escolha evita que uma indisponibilidade momentânea da rede impeça a leitura de uma autorização já concedida, mas também impede que o dispositivo conheça imediatamente uma mudança realizada no servidor. As consequências para revogação, relógio e proteção do segredo serão discutidas na seção dedicada à operação local e à segurança.

A separação também delimita os registros produzidos em cada lado. Solicitação, aprovação e emissão pertencem ao estado central. Leitura, autorização e sinal de acionamento ocorrem no dispositivo. Sem um canal de sincronização e sem sensores adicionais, um registro no banco não equivale automaticamente à abertura da porta, e um comando elétrico não comprova entrada, saída ou ocupação.

\subsection{Decisões tecnológicas}

O frontend foi implementado com Next.js 16 e React 19, usando a organização de rotas pelo diretório \texttt{app}. O App Router do Next.js estrutura páginas e layouts por convenções do sistema de arquivos e integra os componentes React empregados pela interface \cite{vercel2026nextjs}. A escolha permitiu reunir no mesmo projeto as páginas de autenticação, solicitação, consulta e administração, além dos componentes e serviços que consomem a API.

O backend utiliza Python e FastAPI, com modelos Pydantic para representar entradas e saídas das rotas. A documentação oficial do framework destaca a integração entre anotações de tipo, validação de dados e geração de esquemas OpenAPI \cite{fastapi2026features}. No FLIKE, esses recursos organizam os contratos HTTP, enquanto as regras e o acesso ao banco permanecem distribuídos entre rotas, módulos e repositórios próprios do projeto.

A persistência foi implementada com MySQL e o conector oficial para Python. O modelo relacional permite expressar por chaves estrangeiras a hierarquia instituição--edifício--sala--tranca e os vínculos de usuários com solicitações e credenciais. A sintaxe e os tipos empregados seguem o sistema documentado no manual do MySQL \cite{oracle2026mysql84}; a estrutura concreta do FLIKE, entretanto, é definida pelos comandos de criação do próprio backend, apresentados na Subseção~\ref{subsec:modelo-dados}.

No dispositivo físico, a ESP32-CAM reúne câmera e capacidade de processamento para adquirir o QR Code e executar a validação próxima ao atuador. O QR Code foi adotado como transporte visual de dados binários, cuja capacidade depende da versão e do nível de correção de erros \cite{denso2012qressentials}. A autenticação do payload usa AES-CMAC, mecanismo simétrico que permite verificar a integridade e a origem criptográfica da mensagem sob a premissa de proteção do segredo compartilhado \cite{nist2005cmac}. O formato e as consequências de segurança dessas escolhas serão detalhados nas seções dedicadas à credencial digital e à operação local.

\section{Aplicação conectada e persistência}

A aplicação conectada reúne a interface, as regras de negócio e a persistência necessárias para transformar uma solicitação aprovada em credencial. Esta seção parte do modelo de domínio e de dados; as subseções posteriores descreverão a API, o frontend e o fluxo integrado de solicitação e emissão.

\subsection{Modelo de domínio e dados} \label{subsec:modelo-dados}

O esquema relacional foi reconstruído a partir dos comandos SQL de criação executados pelo backend. Essa fonte foi preferida ao arquivo DBML disponível no repositório porque o documento visual omite a tabela de solicitações e conserva uma versão anterior de campos da credencial. O diagrama da Figura~\ref{fig:modelo-dados-flike} representa, portanto, as nove tabelas efetivamente declaradas no script de inicialização e apenas as relações materializadas por chaves estrangeiras.

O núcleo institucional reúne relações sucessivas. \texttt{user} liga-se a \texttt{institution}; a hierarquia segue por \texttt{building} e \texttt{room} até chegar a \texttt{digital\_lock}. A coluna \texttt{owner\_id}, na tabela de instituições, associa cada instituição a um usuário e materializa a responsabilidade contextual: não existe tabela de administradores nem atributo global de papel. A mesma linha de \texttt{user} pode, assim, ser proprietária de instituições e também aparecer como solicitante ou titular de credenciais destinadas a trancas de outra instituição.

A tabela \texttt{digital\_key\_request} registra o usuário solicitante, a tranca destinatária e o estado \texttt{pending}, \texttt{approved} ou \texttt{rejected}. A tabela \texttt{digital\_key} armazena o usuário, a tranca, o payload de 48 bytes, a expiração e os campos de uso previstos pela aplicação. Ambas se ligam independentemente a \texttt{user} e \texttt{digital\_lock}: o esquema não possui uma chave estrangeira entre a solicitação aprovada e a credencial gerada. Essa ausência exige que a correspondência seja inferida pelo fluxo da aplicação e impede demonstrar somente pelo banco uma relação persistente e inequívoca entre os dois registros.

As credenciais de login ficam em \texttt{auth}, separadas dos dados de domínio em \texttt{user}. A associação entre as duas tabelas é conduzida pela aplicação a partir do e-mail transformado, mas não existe \texttt{user\_id} nem chave estrangeira que imponha essa correspondência no banco. Por esse motivo, a Figura~\ref{fig:modelo-dados-flike} usa uma linha tracejada para indicar a associação lógica e a distingue das relações referenciais declaradas pelo esquema.

Por fim, a tabela de eventos, \texttt{event\_log}, prevê registros associados a uma tranca. A tabela \texttt{digital\_key}, por sua vez, possui os campos \texttt{used} e \texttt{used\_at}. A existência dessas estruturas não demonstra coleta automática de eventos do dispositivo: o firmware não mantém comunicação com o backend durante a leitura. Esses campos representam possibilidades do modelo central e não devem ser interpretados como prova de abertura, passagem de uma pessoa ou ocupação do espaço.

\begin{figure}[p]
    \centering
    \includegraphics[width=\textwidth,height=0.76\textheight,keepaspectratio]{modelo_dados_flike.pdf}
    \caption{Modelo de dados efetivo do FLIKE}
    \label{fig:modelo-dados-flike}
    \fonte{Elaborado pelos autores a partir dos comandos SQL do backend.}
\end{figure}
\clearpage

\subsection{Backend}

A API organiza o código em rotas FastAPI, modelos Pydantic, repositórios e infraestrutura de persistência. As rotas recebem os parâmetros HTTP, selecionam o caso de uso e convertem exceções do domínio em respostas. Os modelos delimitam os corpos de entrada e parte das saídas. Os repositórios executam diretamente os comandos SQL de cada entidade, sem uma camada de mapeamento objeto--relacional, enquanto a infraestrutura fornece conexões MySQL obtidas de um pool. As regras de negócio ficam, portanto, distribuídas principalmente entre rotas e repositórios.

O cadastro cria os dados pessoais na tabela \texttt{user} e as credenciais na tabela \texttt{auth}. Nesta última, o e-mail é representado por seu hash e a senha recebe um salt aleatório antes da aplicação de SHA-256. No login, a API repete o cálculo e, se os valores coincidirem, emite um JSON Web Token (JWT). Esse formato transporta um conjunto de declarações entre partes de maneira compacta \cite{jones2015jwt}. Na implementação do FLIKE, o token é protegido com HS256, contém o identificador do usuário e a declaração de expiração, definida para uma hora após sua emissão. O segredo usado nesse mecanismo é uma configuração do servidor distinta do segredo AES-CMAC de cada tranca.

As rotas protegidas recebem o JWT no esquema Bearer, verificam sua integridade e validade e recuperam dele o identificador do usuário. A autorização administrativa não depende de uma classe global de administrador: para cada recurso, a API percorre tranca, sala, edifício e instituição até comparar \texttt{institution.owner\_id} com o usuário autenticado. Esse procedimento delimita consulta e manutenção de instituições, edifícios, salas e trancas próprias, assim como a decisão sobre solicitações destinadas a esses recursos.

As demais rotas podem ser agrupadas por sua função. O catálogo de acesso permite consultar a hierarquia institucional usada pelo solicitante. O conjunto de solicitações cria pedidos e os apresenta separadamente ao solicitante e ao responsável. As operações de decisão rejeitam o pedido ou emitem uma credencial. Por fim, as consultas de credenciais devolvem o payload em representação hexadecimal para que o frontend recupere seus bytes. Essa separação reflete dois contextos do mesmo usuário: solicitar acesso em instituições alheias e administrar a estrutura das instituições que possui.

Uma solicitação autenticada recebe o estado \texttt{pending}, a tranca informada e o identificador extraído do token. Antes de aprová-la ou rejeitá-la, a API confirma que ela permanece pendente e que sua tranca pertence ao responsável autenticado. A rejeição apenas altera o estado para \texttt{rejected}. Na aprovação, o backend cria uma credencial para o solicitante e a tranca e, em seguida, marca o pedido como \texttt{approved}. Quando a requisição não informa outra expiração, o padrão é 24 horas contadas do instante da aprovação e da emissão.

A política do projeto estabelece uma credencial para cada solicitação aprovada, mas essa cardinalidade não é imposta pelo esquema. Como visto na Subseção~\ref{subsec:modelo-dados}, não há chave estrangeira entre pedido e credencial. Além disso, a criação da credencial e a mudança do estado são confirmadas em operações separadas. O fluxo nominal produz a correspondência esperada, porém o banco, isoladamente, não permite reconstruí-la de modo inequívoco nem garante uma transição atômica diante de concorrência ou falha intermediária.

A cobertura de autenticação do commit examinado é parcial. As rotas sob o prefixo administrativo e as operações referentes ao próprio perfil e às próprias solicitações exigem JWT. Em contrapartida, listagens gerais de usuários e da estrutura institucional, consultas de credenciais e operações diretas sobre credenciais não aplicam a mesma dependência. As consultas públicas por usuário ou por identificador devolvem o payload, e as respostas das rotas públicas de trancas incluem \texttt{secret\_key} em hexadecimal.

Essa última exposição rompe uma premissa central do modelo de segurança: quem obtiver o segredo de uma tranca poderá produzir tags AES-CMAC válidas para ela. Assim, a implementação examinada não sustenta uma alegação de proteção efetiva do segredo, ainda que a arquitetura pressuponha seu armazenamento restrito ao backend e ao dispositivo. O fechamento dessas rotas e a exclusão do segredo das respostas são correções necessárias antes de qualquer implantação. As consequências para o modelo de ameaça serão retomadas na seção dedicada à segurança.

\subsection{Frontend}

O frontend utiliza as rotas do App Router para compor as páginas visíveis. Componentes reutilizáveis implementam tabelas, cartões, formulários e modais; hooks coordenam estado, carregamento, mensagens de erro e ações; serviços concentram as requisições HTTP; e tipos TypeScript representam os dados trocados com a API. Essa divisão permite que uma página reúna a apresentação e delegue a comunicação e as transições de estado às demais unidades.

Após o cadastro, o usuário autentica-se e o JWT é guardado pelo navegador. A implementação usa o armazenamento local, denominado \textit{localStorage} pela API Web, sob a chave \texttt{access\_token}. As páginas conectadas verificam sua presença para redirecionar ao login, mas essa verificação local não valida o token nem constitui autorização. A decisão efetiva ocorre quando uma rota protegida da API valida o Bearer recebido. Não há, no código examinado, middleware global de sessão, renovação do token ou persistência do login fora do armazenamento do navegador.

No fluxo de solicitação, a interface conduz o usuário pela seleção sucessiva de instituição, edifício e sala. Depois da escolha da sala, o hook consulta as trancas associadas e utiliza a primeira retornada pela API; a versão examinada não apresenta ao usuário uma escolha quando há múltiplas trancas. A ação final envia apenas o identificador dessa tranca, pois o backend deriva do JWT a identidade do solicitante. Uma confirmação textual informa que o pedido foi registrado.

O painel do solicitante combina os dados do usuário, suas credenciais e os pedidos rejeitados. Para apresentar uma credencial, o frontend consulta o registro pelo identificador, separa os pares hexadecimais, reconstrói um vetor de bytes e o entrega à biblioteca de QR Code em modo binário. O resultado é uma URL de dados exibida em um modal. Essa geração ocorre no navegador e não cria nova credencial no servidor. Também não implementa, por si só, download dedicado, funcionamento como aplicativo móvel ou persistência da imagem depois de uma recarga.

A classificação exibida no painel ainda considera ativa apenas a credencial não marcada como usada e cuja expiração está no futuro. Esse critério conserva a semântica legada do campo \texttt{used}, mas não representa a política final do FLIKE: a mesma credencial pode ser reapresentada durante sua janela de validade. Por isso, a marcação de uso da aplicação conectada não será empregada como regra de autorização física.

No fluxo do responsável, a mesma aplicação apresenta somente instituições, edifícios, salas e trancas que pertencem ao usuário autenticado. A tela de pedidos permite filtrar os estados \texttt{pending}, \texttt{approved} e \texttt{rejected}. Os botões de aprovação e rejeição aparecem para pedidos pendentes e, depois da resposta da API, a listagem é consultada novamente. Embora os caminhos e rótulos da interface usem o termo administrativo, a permissão continua derivada da propriedade da instituição.

\subsection{Fluxo de solicitação e emissão}

A Figura~\ref{fig:fluxo-solicitacao-emissao} reúne o caminho principal sem antecipar o formato interno da credencial. Na primeira etapa, o solicitante escolhe a tranca por meio da aplicação. O frontend envia o pedido autenticado, a API associa o usuário extraído do token e o banco persiste o estado pendente. A resposta confirma o registro, mas ainda não contém uma autorização.

Na segunda etapa, o responsável consulta os pedidos relacionados às instituições que possui. Quando envia sua decisão, a API verifica novamente a propriedade da tranca e o estado do pedido. Uma rejeição encerra o fluxo administrativo sem credencial. Uma aprovação cria o payload autenticado, persiste a credencial e altera o pedido para \texttt{approved}. Por padrão, sua janela termina 24 horas depois desse instante; uma expiração diferente pode ser informada no contrato da API.

\begin{figure}[htbp]
    \centering
    \includegraphics[width=\textwidth,height=0.78\textheight,keepaspectratio]{fluxo_solicitacao_emissao.pdf}
    \caption{Sequência de solicitação, decisão e emissão da credencial}
    \label{fig:fluxo-solicitacao-emissao}
    \fonte{Elaborado pelos autores a partir da implementação do FLIKE.}
\end{figure}

Em consulta posterior, o painel do solicitante recupera a credencial emitida e renderiza o QR Code. A obtenção do payload depende da aplicação conectada; somente a leitura e a decisão junto à fechadura dispensam consulta ao servidor. A aprovação, a emissão, a obtenção da imagem e a apresentação física são eventos distintos. Uma única emissão nominal decorre da solicitação aprovada, enquanto a credencial resultante permanece reapresentável até o instante de expiração, exclusive.

\section{Credencial digital}

\subsection{Conceito e ciclo de vida}

A credencial digital é a representação portável de uma autorização temporária. Ela nasce quando o responsável aprova uma solicitação, permanece associada ao usuário lógico e a uma tranca específica e deixa de ser aceita ao término do intervalo concedido. Solicitação, credencial e QR Code têm papéis distintos: a solicitação registra o pedido e sua decisão; a credencial reúne os dados autenticados da autorização; e o QR Code é somente o meio óptico pelo qual os bytes chegam ao dispositivo.

O segredo criptográfico também não integra a credencial. O backend e a tranca destinatária possuem previamente uma cópia desse valor, provisionado no dispositivo pelo fornecedor na programação da ESP32-CAM. O servidor usa o segredo para produzir o código de autenticação, e o dispositivo o usa para verificar o código recebido. Esse arranjo permite a decisão local, mas faz da proteção do segredo uma premissa de segurança do sistema.

O ciclo nominal possui quatro transições. Primeiro, o servidor registra a solicitação pendente. Segundo, o responsável a aprova ou rejeita. Terceiro, somente uma aprovação produz a credencial, cuja emissão inicia a janela de autorização. Quarto, o titular apresenta o QR Code à tranca durante essa janela. A apresentação não consome a autorização: o mesmo símbolo pode ser lido novamente até o instante de expiração.

\subsection{Informações autenticadas}

O conteúdo protegido reúne quatro valores. \texttt{user\_id} identifica a conta à qual a autorização foi concedida; \texttt{digital\_lock\_id} identifica a tranca lógica destinatária; \texttt{issued\_at} registra o primeiro instante admissível; e \texttt{expires\_at} marca o limite superior, que já não pertence ao intervalo. O usuário físico não é autenticado pela apresentação: quem obtiver uma cópia do QR Code poderá reapresentá-lo enquanto ele for válido.

Os instantes são representados como segundos inteiros desde a época POSIX. A condição temporal adotada é

\begin{equation}
    \texttt{issued\_at} \leq \texttt{agora} < \texttt{expires\_at}.
\end{equation}

Essa definição evita ambiguidade na fronteira: a credencial vale no instante de emissão, inclusive, e deixa de valer no instante de expiração, exclusive. A comparação depende de um relógio coerente no dispositivo; operar sem servidor não elimina essa dependência temporal.

\subsection{Serialização e autenticação}

Antes do cálculo criptográfico, cada um dos quatro valores é convertido em um inteiro sem sinal de 64 bits, em ordem de bytes do mais significativo para o menos significativo (\textit{big-endian}). A concatenação produz 32 bytes de dados. Sobre essa sequência, o backend calcula uma tag AES-CMAC de 16 bytes com uma chave AES de 256 bits e a anexa ao final. O payload completo possui, como consequência, 48 bytes. O Quadro~\ref{qua:layout-credencial} apresenta o leiaute.

\begin{quadro}[htb]
    \caption{Leiaute binário da credencial do FLIKE}
    \label{qua:layout-credencial}
    \small
    \setlength{\tabcolsep}{4pt}
    \begin{tabular}{|p{0.20\textwidth}|p{0.12\textwidth}|p{0.18\textwidth}|p{0.40\textwidth}|}
        \hline
        \textbf{Campo} & \textbf{Bytes} & \textbf{Intervalo} & \textbf{Função} \\ \hline
        \texttt{user\_id} & 8 & 0--7 & Identificar a conta autorizada. \\ \hline
        \texttt{digital\_lock\_id} & 8 & 8--15 & Vincular a autorização à tranca destinatária. \\ \hline
        \texttt{issued\_at} & 8 & 16--23 & Delimitar o início da autorização em tempo POSIX. \\ \hline
        \texttt{expires\_at} & 8 & 24--31 & Delimitar o fim exclusivo da autorização. \\ \hline
        tag AES-CMAC & 16 & 32--47 & Autenticar conjuntamente os 32 bytes anteriores. \\ \hline
    \end{tabular}
    \fonte{Elaborado pelos autores a partir da implementação do backend.}
\end{quadro}

O CMAC é um código de autenticação de mensagem baseado em cifra de bloco. Sob a premissa de que somente os componentes autorizados conhecem a chave, a coincidência da tag fornece evidência de que os dados foram produzidos por alguém que possui esse segredo e não foram alterados sem a recomputação do código \cite{nist2005cmac}. O mecanismo não oculta os identificadores nem os instantes, não impede a cópia de uma mensagem autêntica e não produz assinatura digital ou não repúdio. A comparação deve abranger a tag completa, e qualquer diferença deve resultar em rejeição.

\subsection{Codificação em QR Code}

O frontend recebe o payload em hexadecimal, converte cada par para o byte correspondente e entrega o vetor à biblioteca de QR Code em modo binário. Essa etapa é necessária porque tratar o conteúdo como texto introduziria uma nova codificação e poderia alterar bytes nulos ou valores superiores a \texttt{0x7f}. O símbolo final é fixado na versão 3, cuja matriz possui 29 por 29 módulos, e usa nível L de correção de erros.

A escolha decorreu de um limite observado durante o desenvolvimento: símbolos mais densos eram lidos com menor facilidade pela câmera do protótipo. O formato binário compacto possibilitou transportar os 48 bytes na versão 3-L, mantendo a matriz menor. O nível L reserva menos redundância do que níveis superiores e, portanto, aceita menor dano físico; trata-se de um compromisso entre capacidade, densidade óptica e tolerância a degradação, não de uma propriedade criptográfica \cite{denso2012qressentials}. Além disso, o QR Code não é um canal confiável por natureza: ele pode ser copiado ou substituído, razão pela qual a decisão depende do conteúdo autenticado e do destino conferido pelo firmware \cite{kieseberg2010qrcode}.

\subsection{Decisão de aceitação}

A leitura bem-sucedida inicia a validação, mas não autoriza o acesso. O dispositivo confirma o comprimento completo, separa dados e tag, recalcula o AES-CMAC, compara o identificador destinatário e avalia a janela temporal. Somente a conjunção dessas condições produz o sinal lógico de abertura. A Figura~\ref{fig:validacao-local-flike} representa o fluxo; sua ordem gráfica não altera a exigência de que todas as verificações precedam o acionamento.

\begin{figure}[htbp]
    \centering
    \includegraphics[width=0.90\textwidth,height=0.78\textheight,keepaspectratio]{validacao_local_flike.pdf}
    \caption{Fluxo local de validação da credencial}
    \label{fig:validacao-local-flike}
    \fonte{Elaborado pelos autores.}
\end{figure}

Uma tag válida para outra tranca continua sendo uma mensagem autêntica, mas não uma autorização para o atuador que a recebeu. De forma semelhante, uma mensagem autêntica e destinada ao dispositivo ainda deve ser recusada fora do intervalo temporal. O estado \texttt{used} mantido no banco não participa dessa decisão, pois o dispositivo deliberadamente não consulta o servidor e a política final permite reapresentações durante a vigência.

\section{Dispositivo físico}

\subsection{Plataforma embarcada e aquisição do QR Code}

O dispositivo foi construído com uma placa ESP32-CAM no arranjo AI-Thinker e uma câmera OV2640. A biblioteca de leitura inicializa a câmera, executa a aquisição em uma tarefa do sistema operacional de tempo real e disponibiliza o vetor bruto e seu comprimento ao laço principal. O processamento preserva o tamanho fornecido pelo decodificador; ele não interpreta a credencial como cadeia de caracteres terminada por zero.

A ESP32 integra processador, memória, periféricos digitais e recursos de comunicação adequados à prototipação de sistemas embarcados. No FLIKE, sua função relevante é concentrar aquisição óptica, validação da credencial e geração do sinal de controle próximo ao atuador. O produto não depende de seus rádios durante a decisão de acesso. A ausência de tráfego nesse instante reduz a dependência operacional, mas não transforma a placa em componente endurecido contra extração física de segredo.

\subsection{Firmware de validação}

O firmware preservado no repositório registra a integração entre a biblioteca da câmera e o cálculo AES-CMAC da biblioteca Mbed TLS. Ele separa os 32 bytes de dados dos 16 bytes finais, recalcula a tag com AES-256 e compara o resultado. Essa versão arquivada imprime o conteúdo e o resultado no monitor serial; não contém, no mesmo commit, todas as verificações e o acionamento usados no ensaio final.

Na versão gravada para a demonstração integrada, o processamento também confirmou o comprimento, recuperou os quatro inteiros em ordem \textit{big-endian}, comparou \texttt{digital\_lock\_id} com o identificador provisionado, avaliou \texttt{issued\_at} e \texttt{expires\_at} e emitiu \texttt{HIGH} somente após a aceitação. O código-fonte exato desse ensaio não foi preservado no repositório. Essa lacuna impede reproduzir a imagem de firmware a partir do controle de versão, mas não elimina o resultado observado na demonstração física, descrito na Subseção~\ref{subsec:demonstracao-fisica}.

\subsection{Circuito de acionamento}

O estágio elétrico reaproveita a topologia construída pelos autores em um projeto anterior da disciplina Laboratório de Processadores, antecedente histórico do FLIKE. O sinal lógico da ESP32-CAM chega à base de um transistor NPN 2N2222 por um resistor de 1~k$\Omega$. O transistor funciona como chave e condiciona o comando aplicado à entrada do módulo de relé de 12~V; um segundo resistor de 1~k$\Omega$ estabelece o nível de repouso. O contato normalmente aberto do relé comuta a alimentação da fechadura Papaiz AA-ERL200P.

O circuito usa duas alimentações. Um carregador de celular adaptado fornece a baixa tensão da ESP32-CAM, enquanto uma fonte chaveada de 12~V alimenta relé e fechadura. Essa separação permite que o pino digital controle uma carga que não pode ser acionada diretamente pelo microcontrolador. A Figura~\ref{fig:circuito-eletrico-flike} registra a topologia; o ramo de acionamento manual representado no esquema foi previsto, mas não montado no protótipo avaliado.

\begin{figure}[p]
    \centering
    \includegraphics[width=\textwidth,height=0.78\textheight,keepaspectratio]{circuito_eletrico_flike.png}
    \caption{Esquema elétrico do estágio de acionamento empregado no FLIKE}
    \label{fig:circuito-eletrico-flike}
    \fonte{Elaborado pelos autores.}
\end{figure}
\clearpage

O protótipo é uma montagem de bancada em protoboard, sem placa dedicada, gabinete, bateria ou certificação elétrica. O relé oferece separação funcional entre comando e carga, mas o projeto não avaliou isolação, imunidade eletromagnética, aquecimento, durabilidade ou adequação a uma instalação permanente. A Figura~\ref{fig:prototipo-flike} identifica os componentes principais da montagem.

\begin{figure}[htbp]
    \centering
    \includegraphics[width=\textwidth,height=0.72\textheight,keepaspectratio]{prototipo_flike_anotado.png}
    \caption{Protótipo experimental do circuito de acionamento}
    \label{fig:prototipo-flike}
    \fonte{Acervo dos autores.}
\end{figure}

\subsection{Demonstração física} \label{subsec:demonstracao-fisica}

A equipe realizou uma demonstração física ponta a ponta do FLIKE em 31 de agosto de 2026. Uma credencial foi apresentada à câmera; o dispositivo reconheceu o QR Code, recuperou os 48 bytes, extraiu os campos e verificou comprimento, tranca destinatária, instante de emissão, expiração e tag AES-CMAC. Diante do caso válido, o microcontrolador emitiu o nível \texttt{HIGH}; o transistor comandou o relé; e a fechadura elétrica foi acionada.

O ensaio estabeleceu a continuidade funcional entre canal óptico, decisão embarcada e carga elétrica. Seu registro disponível é o relato concordante da equipe, apoiado pelo código preservado dos subsistemas, pelo esquema e pela fotografia da montagem. Não foram preservados vídeo, log serial ou imagem binária do firmware do ensaio. Consequentemente, o resultado sustenta que a integração ocorreu, mas não permite calcular taxa de sucesso, tempo médio de leitura, tolerância a ângulo ou iluminação, consumo, repetibilidade estatística ou confiabilidade de longo prazo.

O evento observado foi o acionamento da fechadura. Como não há sensor de estado da porta, de passagem ou de ocupação, o sistema não distingue automaticamente destravamento, abertura física, entrada, saída e permanência no ambiente.

\section{Operação local, segurança e auditoria}

\subsection{Consequências da operação sem servidor}

A decisão local permite utilizar uma credencial já emitida quando o dispositivo não possui acesso à API. Essa propriedade cobre leitura, validação e acionamento; cadastro, solicitação, aprovação, emissão e recuperação inicial do QR Code continuam conectados. Ela exige que identificador, segredo e relógio estejam disponíveis localmente.

O mesmo isolamento impede revogação imediata. Depois da emissão, retirar a permissão no banco não informa a tranca sem um canal adicional; a credencial poderá ser aceita até \texttt{expires\_at}. Janelas curtas reduzem a duração desse risco, mas aumentam a frequência de emissão e a dependência do serviço web. O padrão de 24 horas adotado no backend é uma política do protótipo e deve ser ajustado ao risco do espaço em uma implantação.

A cópia do símbolo é outro efeito relevante. Como o dispositivo não consulta um estado de consumo, qualquer cópia byte a byte permanece equivalente ao original durante a vigência. A política do FLIKE admite reapresentação, mas não pretende autorizar o compartilhamento. Vincular o identificador da conta aos dados oferece rastreabilidade lógica; não autentica fisicamente quem segura o telefone ou a imagem impressa.

\subsection{Proteção do segredo e modelo de ameaça}

O modelo considera um adversário capaz de observar, copiar e alterar o QR Code e de chamar interfaces web expostas, mas pressupõe que não conheça o segredo da tranca. Sob essa condição, alterações dos quatro campos são detectadas pelo AES-CMAC. A proteção do material criptográfico durante geração, armazenamento, distribuição e substituição é condição necessária para manter essa propriedade \cite{barker2020keymanagement}.

No protótipo, o fornecedor grava o identificador e o segredo ao programar o dispositivo. Não há interface de provisionamento, rotação, revogação de chave criptográfica nem elemento seguro. A extração do segredo do servidor ou da placa permite fabricar credenciais aceitas para aquela tranca. Um mesmo segredo compartilhado por várias trancas simplifica salas com mais de um atuador, mas amplia o impacto de um comprometimento.

A inspeção da API identificou uma violação direta dessa premissa: respostas públicas de consulta de trancas incluem \texttt{secret\_key} em hexadecimal. Também há consultas públicas de credenciais que expõem payloads. Antes de qualquer uso real, essas rotas precisam exigir autenticação e autorização contextual, e o segredo precisa ser removido dos esquemas de resposta e de registros operacionais. Enquanto essa correção não estiver presente, o protótipo não oferece proteção efetiva contra fabricação de credenciais por um cliente da API.

O AES-CMAC fornece autenticação simétrica e integridade, não irretratabilidade. Backend e dispositivo compartilham o mesmo segredo e ambos podem produzir tags indistinguíveis. A atribuição de uma solicitação a uma conta e de uma apresentação a uma credencial depende conjuntamente de autenticação web, proteção do dispositivo e confiabilidade dos registros; mesmo assim, não comprova que a pessoa física titular executou a ação.

\subsection{Eventos e limites da auditoria}

Para interpretar os registros, o FLIKE distingue seis eventos: emissão da credencial pelo servidor; leitura do símbolo pela câmera; aceitação ou rejeição lógica; emissão do sinal de controle; movimento da fechadura; e passagem de uma pessoa. O banco registra pedidos e credenciais emitidas. A implementação local decide e aciona, porém não sincroniza um log embarcado com o servidor. A fechadura não fornece realimentação ao microcontrolador, e não há sensor de porta ou de presença.

Assim, o histórico administrativo comprova autorizações concedidas no domínio da aplicação. Ele não constitui histórico de acessos físicos. O campo \texttt{used} e a tabela \texttt{event\_log} previstos no esquema não mudam essa conclusão, pois o dispositivo não os atualiza durante a operação offline. Uma auditoria física exigiria armazenamento confiável no dispositivo, sincronização autenticada e sensores compatíveis com o evento que se pretende afirmar.

\section{Verificação técnica e resultados}

\subsection{Estratégia e ambiente}

A verificação combinou três fontes. A inspeção estática examinou os commits finais indicados pela equipe para backend e frontend, o commit preservado do firmware, o esquema SQL, as dependências e os artefatos elétricos. Uma verificação executável independente reproduziu a serialização e o AES-CMAC com a biblioteca \texttt{cryptography}. Por fim, as demonstrações de software e hardware confirmadas pela equipe foram usadas como evidência de integração. Nenhuma dessas fontes foi apresentada como teste com usuários.

Para o vetor criptográfico, adotaram-se uma chave de 32 bytes com valores sucessivos de 0 a 31, usuário 17, tranca 23, emissão 1770000000 e expiração 1770086400. A serialização produziu os 32 bytes esperados, e o AES-CMAC produziu a tag \texttt{2e98 1020 9a16 6e7e 1a61 ece3 2565 7c45}. A alteração de um bit em cada um dos quatro campos, sem recomputar a tag, foi rejeitada nos quatro casos. O ensaio verifica a regra criptográfica e a compatibilidade do formato; ele não avalia resistência da chave nem o armazenamento do segredo.

Não foram executadas medições sistemáticas de leitura óptica, ensaios ambientais, testes elétricos instrumentados ou avaliação de usabilidade. Os fluxos de usuário e responsável e o acionamento físico são resultados de demonstração, sem amostra repetida que permita inferência estatística.

\subsection{Resultados por subsistema}

Na aplicação conectada, a demonstração percorreu cadastro e autenticação, solicitação de acesso, cadastro administrativo de edifícios e trancas, aprovação ou rejeição e disponibilização da credencial no painel do solicitante. A inspeção confirmou os contratos correspondentes nos dois repositórios. Também encontrou lacunas: a interface seleciona implicitamente a primeira tranca de uma sala, a ligação entre pedido e credencial não é persistida por chave estrangeira, e rotas sensíveis permanecem públicas.

Na credencial, a implementação e o vetor executado confirmaram a ordem dos campos, o tamanho, a ordem de bytes e a detecção de alterações pelo AES-CMAC. Os testes de QR Code realizados durante o desenvolvimento confirmaram o transporte dos bytes e levaram à fixação da versão 3-L. Não foi preservada uma série de tempos de leitura; portanto, não se apresenta valor de desempenho.

No dispositivo físico, a demonstração final confirmou leitura, decisão local e acionamento sem consulta ao servidor. A montagem elétrica respondeu ao sinal \texttt{HIGH} e destravou a fechadura. O resultado atende ao objetivo de prototipação ponta a ponta, com as limitações de registro e repetibilidade já expostas.

\subsection{Matriz de atendimento dos requisitos}

Os estados dos Quadros~\ref{qua:resultados-rf-web} a~\ref{qua:resultados-ra} usam quatro categorias. \textit{Atendido} indica evidência suficiente para o critério no nível do protótipo; \textit{parcial} indica que parte do comportamento existe, mas uma condição falhou ou não foi coberta; \textit{não atendido} indica conflito direto com a implementação; e \textit{não avaliado} indica ausência de evidência para decidir.

\begin{quadro}[htbp]
    \caption{Resultados de RF-01 a RF-08}
    \label{qua:resultados-rf-web}
    \small
    \begin{tabular}{|p{0.10\textwidth}|p{0.15\textwidth}|p{0.61\textwidth}|}
        \hline
        \textbf{ID} & \textbf{Estado} & \textbf{Evidência e ressalva principal} \\ \hline
        RF-01 & Atendido & Cadastro, autenticação, JWT e rejeição de credencial incorreta. \\ \hline
        RF-02 & Atendido & Propriedade contextual e atuação como solicitante e responsável. \\ \hline
        RF-03 & Parcial & Hierarquia existe; autorização e exclusões não foram cobertas em todas as operações. \\ \hline
        RF-04 & Não atendido & Interface escolhe a primeira tranca da sala sem seleção explícita. \\ \hline
        RF-05 & Atendido & Criação e acompanhamento nos três estados. \\ \hline
        RF-06 & Atendido & Decisão pelo proprietário e transição de estado. \\ \hline
        RF-07 & Parcial & Aprovação emite; banco não garante vínculo único nem atomicidade. \\ \hline
        RF-08 & Parcial & QR é apresentado; consultas públicas impedem comprovar isolamento. \\ \hline
    \end{tabular}
    \fonte{Elaborado pelos autores.}
\end{quadro}

\begin{quadro}[htbp]
    \caption{Resultados de RF-09 a RF-14}
    \label{qua:resultados-rf-credencial}
    \small
    \begin{tabular}{|p{0.10\textwidth}|p{0.15\textwidth}|p{0.61\textwidth}|}
        \hline
        \textbf{ID} & \textbf{Estado} & \textbf{Evidência e ressalva principal} \\ \hline
        RF-09 & Atendido & Quatro valores da autorização são serializados e autenticados. \\ \hline
        RF-10 & Atendido & Vetor executado e alterações individuais rejeitadas. \\ \hline
        RF-11 & Atendido & Transporte e recuperação dos 48 bytes demonstrados. \\ \hline
        RF-12 & Atendido & Estrutura, tag, destino e intervalo conferidos antes do acionamento. \\ \hline
        RF-13 & Atendido & Reapresentação aceita na janela e rejeitada fora dela. \\ \hline
        RF-14 & Parcial & Consultas existem, mas sua proteção é insuficiente. \\ \hline
    \end{tabular}
    \fonte{Elaborado pelos autores.}
\end{quadro}

\begin{quadro}[htbp]
    \caption{Resultados dos requisitos não funcionais}
    \label{qua:resultados-rnf}
    \small
    \begin{tabular}{|p{0.10\textwidth}|p{0.15\textwidth}|p{0.61\textwidth}|}
        \hline
        \textbf{ID} & \textbf{Estado} & \textbf{Evidência e ressalva principal} \\ \hline
        RNF-01 & Atendido & Leitura, decisão e acionamento sem consulta à API. \\ \hline
        RNF-02 & Parcial & Operações principais verificam propriedade; há rotas sensíveis públicas. \\ \hline
        RNF-03 & Não atendido & API pública expõe o segredo em hexadecimal. \\ \hline
        RNF-04 & Não avaliado & Roteiro completo não foi executado em viewports definidos. \\ \hline
        RNF-05 & Atendido & QR Code 3-L e leitura pela ESP32-CAM. \\ \hline
    \end{tabular}
    \fonte{Elaborado pelos autores.}
\end{quadro}

\begin{quadro}[htbp]
    \caption{Resultados dos requisitos de acessibilidade}
    \label{qua:resultados-ra}
    \small
    \begin{tabular}{|p{0.10\textwidth}|p{0.15\textwidth}|p{0.61\textwidth}|}
        \hline
        \textbf{ID} & \textbf{Estado} & \textbf{Evidência e ressalva principal} \\ \hline
        RA-01 & Atendido & Apresentação e destravamento sem funcionário. \\ \hline
        RA-02 & Parcial & Há mensagens textuais; inspeção de todos os estados não foi executada. \\ \hline
        RA-03 & Parcial & Sequência é explícita, exceto escolha entre múltiplas trancas. \\ \hline
        RA-04 & Atendido & Sem alarme, música ou luz intermitente como feedback. \\ \hline
    \end{tabular}
    \fonte{Elaborado pelos autores.}
\end{quadro}
\clearpage

Dos 23 requisitos avaliados, 13 foram classificados como atendidos, sete como parciais, dois como não atendidos e um como não avaliado. A contagem não constitui uma nota global: a exposição do segredo, por exemplo, tem impacto maior sobre a segurança do que a ausência de uma medição de responsividade. O resultado central é a viabilidade funcional da cadeia ponta a ponta, acompanhada por lacunas que impedem tratar o protótipo como produto pronto para implantação.

\section{Limitações da implementação}

A primeira limitação é metodológica. Os requisitos de acessibilidade foram definidos pela equipe e fundamentados por literatura, porém não foram validados com pessoas autistas ou outros usuários da sala. O projeto pode reduzir uma interação presencial obrigatória no fluxo concebido, mas não demonstrou redução de constrangimento, ansiedade, estímulos sensoriais ou esforço cognitivo.

A segunda é de segurança operacional. A decisão offline não recebe revogação imediata, depende do relógio local e aceita cópias da credencial durante a validade. O segredo é provisionado manualmente, sem rotação nem hardware seguro, e a versão inspecionada da API o expõe. Essas condições exigem correção antes de qualquer piloto com controle de acesso real.

A terceira é de rastreabilidade física. O sistema não possui sensor de porta, identificação da pessoa que apresenta o QR Code, registro embarcado persistente ou sincronização de eventos. Por isso, emissão, leitura, autorização, acionamento, abertura e passagem não podem ser confundidos. Os registros centrais descrevem permissões concedidas, não ocupação do ambiente.

A quarta é de maturidade da engenharia. O firmware exato da demonstração não foi preservado, faltam testes automatizados entre Python e C++, e a montagem permanece em protoboard com fontes de bancada. Não foram avaliados falhas de energia, durabilidade, temperatura, ataque físico, desempenho óptico ou conformidade elétrica. Também persistem inconsistências de software, como seleção implícita da primeira tranca, campos legados de uso, transações não atômicas e rotas sem proteção adequada.

Apesar dessas limitações, o desenvolvimento materializou o percurso completo proposto: uma decisão administrativa no sistema web originou uma credencial compacta, o dispositivo verificou essa autorização localmente e o circuito acionou a fechadura. O capítulo seguinte confronta esse resultado com a pergunta e os objetivos do trabalho.
